\subsection{What Is an Ontology}
\label{sec:what-is-an-ontology}
\slidesrc{Ontologies: 3--7, 10--21}

\subsubsection{Two definitions}
\slidesrc{Ontologies: 3--7}

The word comes from philosophy, where \emph{ontology} is the study of being. Computer
science borrowed it in the early 1990s with a much more practical meaning:
\begin{itemize}[nosep]
  \item ``The study of that which is common to all things which exist''
        (Aristotle, \emph{Metaphysics});
  \item ``\textbf{Explicit specification of a conceptualization}'' (Tom Gruber,
        \emph{A Translation Approach to Portable Ontology Specifications}, Knowledge
        Acquisition 5(2):199--220, 1993).
\end{itemize}

\begin{nbintuition}[reading Gruber's definition]
A \emph{conceptualization} is the way we carve up a domain: which kinds of things exist
(wines, producers, foods), how they are related (a wine \emph{is produced by} a
producer), which constraints hold. It normally lives only in people's heads.
\emph{Explicit specification} means writing it down, in a formal language, so that it
can be shared between people and programs and checked by a machine.
\end{nbintuition}

\subsubsection{Motivation: semantic interoperability}
\slidesrc{Ontologies: 10--12}

\begin{nbdef}{Semantic interoperability}
\textbf{Semantic interoperability} is ``the ability of computer systems to exchange data
with unambiguous, shared meaning''.
\end{nbdef}

Two systems can exchange bytes perfectly and still disagree on what they mean: is
\texttt{price} with or without VAT, is \texttt{date} the order date or the shipping
date? Semantic interoperability has been one of the most active research fields of the
last twenty-five years: the advent of ontologies in the early nineties, Web Services at
the beginning of the new millennium, and their combination into Semantic Web Services
all aim at making the ``semantic interoperability dream'' come true.

\begin{nbwhy}
The interest is economic as much as scientific. Walker et al.\ (\emph{The Value of
Healthcare Information Exchange and Interoperability}, Health Affairs, 2005) estimated
that a fully standardised system for information exchange and interoperability in US
health care could yield a net value of \textbf{77.8 billion dollars per year} once
implemented. More recent work (e.g.\ on semantic interoperability in health care, 2018,
and the European Data Conference on Reference Data and Semantics, 2023) confirms the
expected advantages.
\end{nbwhy}

\subsubsection{Gruber's definition, 2008}
\slidesrc{Ontologies: 13--15}

Fifteen years after his first definition, Gruber gave a more detailed one (\emph{Ontology},
Encyclopedia of Database Systems, Springer, 2008).

\begin{nbdef}{Ontology (Gruber, 2008)}
In computer and information science, an \textbf{ontology} defines a set of
\textbf{representational primitives} with which to model a domain of knowledge or
discourse. The primitives are typically
\begin{itemize}[nosep]
  \item \textbf{classes} (or sets),
  \item \textbf{attributes} (or properties),
  \item \textbf{relationships} (relations among class members).
\end{itemize}
Their definitions include information about their \emph{meaning} and \emph{constraints}
on their logically consistent application.
\end{nbdef}

From the point of view of databases:
\begin{itemize}[nosep]
  \item an ontology is a level of abstraction of data models, analogous to hierarchical
        and relational models, but intended for modelling \emph{knowledge} about
        individuals, their attributes and their relationships to other individuals;
  \item ontology languages abstract away from data structures and implementation
        strategies, and are closer in expressive power to \textbf{first-order logic} than
        database languages. Ontologies are therefore said to be at the
        \textbf{``semantic'' level}, database schemas at the ``logical'' or ``physical''
        level;
  \item being independent of the lower-level data models, ontologies are used to
        integrate heterogeneous databases, to make disparate systems interoperate and to
        specify interfaces to knowledge-based services. In the Semantic Web standards
        they are an explicit layer (Section~\ref{sec:semantic-web}), with standard
        languages and many commercial and open source tools.
\end{itemize}

\subsubsection{An example: a wine ontology}
\slidesrc{Ontologies: 16}

The slides show a small ontology of wines and foods (in Italian). Each node is a class;
each child is a subclass (``is a kind of'') of its parent, exactly like the hypernym
hierarchy of WordNet (Section~\ref{sec:wordnet}).

\begin{center}
\begin{tikzpicture}[
    font=\small, x=0.6cm, y=-0.48cm,
    cls/.style={anchor=west, inner sep=2pt},
    inner/.style={cls, text=nbblue!80!black},
    edge/.style={draw=nbgray, thick}]
  \node[inner] (thing) at (0,0) {Thing};
  \node[inner] (vino) at (1,1) {Vino \emph{(wine)}};
  \node[inner] (vb) at (2,2) {Vino Bianco \emph{(white wine)}};
  \node[cls] (pb) at (3,3) {Pinot Bianco};
  \node[inner] (ri) at (3,4) {Riesling};
  \node[cls] (rs) at (4,5) {Riesling Secco \emph{(dry)}};
  \node[cls] (rd) at (4,6) {Riesling Dolce \emph{(sweet)}};
  \node[cls] (ch) at (3,7) {Chardonnay};
  \node[cls] (vr) at (2,8) {Vino Rosso \emph{(red)}};
  \node[cls] (vro) at (2,9) {Vino Rosè \emph{(rosé)}};
  \node[cls] (vd) at (2,10) {Vino da dessert \emph{(dessert wine)}};
  \node[inner] (cibo) at (1,11) {Cibo \emph{(food)}};
  \node[inner] (carne) at (2,12) {Carne \emph{(meat)}};
  \node[cls] (cb) at (3,13) {Carni Bianche \emph{(white meat)}};
  \node[cls] (cr) at (3,14) {Carni Rosse \emph{(red meat)}};
  \node[inner] (pesce) at (2,15) {Pesce \emph{(fish)}};
  \node[cls] (ps) at (3,16) {Pesce d'acqua salata \emph{(saltwater)}};
  \node[cls] (pd) at (3,17) {Pesce d'acqua dolce \emph{(freshwater)}};
  \node[cls] (prod) at (1,18) {Produttore \emph{(producer)}};
  \draw[edge] ([xshift=0.3cm]thing.south west) |- (vino.west);
  \draw[edge] ([xshift=0.3cm]vino.south west) |- (vb.west);
  \draw[edge] ([xshift=0.3cm]vb.south west) |- (pb.west);
  \draw[edge] ([xshift=0.3cm]vb.south west) |- (ri.west);
  \draw[edge] ([xshift=0.3cm]ri.south west) |- (rs.west);
  \draw[edge] ([xshift=0.3cm]ri.south west) |- (rd.west);
  \draw[edge] ([xshift=0.3cm]vb.south west) |- (ch.west);
  \draw[edge] ([xshift=0.3cm]vino.south west) |- (vr.west);
  \draw[edge] ([xshift=0.3cm]vino.south west) |- (vro.west);
  \draw[edge] ([xshift=0.3cm]vino.south west) |- (vd.west);
  \draw[edge] ([xshift=0.3cm]thing.south west) |- (cibo.west);
  \draw[edge] ([xshift=0.3cm]cibo.south west) |- (carne.west);
  \draw[edge] ([xshift=0.3cm]carne.south west) |- (cb.west);
  \draw[edge] ([xshift=0.3cm]carne.south west) |- (cr.west);
  \draw[edge] ([xshift=0.3cm]cibo.south west) |- (pesce.west);
  \draw[edge] ([xshift=0.3cm]pesce.south west) |- (ps.west);
  \draw[edge] ([xshift=0.3cm]pesce.south west) |- (pd.west);
  \draw[edge] ([xshift=0.3cm]thing.south west) |- (prod.west);
\end{tikzpicture}
\captionof{figure}{The wine ontology of the slides as a class hierarchy (translations in
italics).}
\label{fig:wine-ontology}
\end{center}

The slide also labels the \textbf{levels} of the hierarchy: \emph{top level} (the most
general classes, such as \emph{Vino}), \emph{middle level} (\emph{Vino Bianco},
\emph{Riesling}, \emph{Chardonnay}, \ldots) and \emph{bottom level} (the most specific
classes, \emph{Riesling Secco} and \emph{Riesling Dolce}). A full ontology would also
have \emph{properties} linking the hierarchies, e.g.\ a wine \emph{is made by} a
\emph{Produttore}, a wine \emph{goes well with} a \emph{Cibo}.

\subsubsection{Ontologies and databases}
\slidesrc{Ontologies: 17--18}

\begin{center}
\small
\begin{tabularx}{\linewidth}{@{}l >{\raggedright\arraybackslash}X >{\raggedright\arraybackslash}X@{}}
\toprule
 & \textbf{Ontologies} & \textbf{Databases} \\
\midrule
Purpose & sharing, reuse, defining semantics & storing large amounts of structured
  data \\
Expressive power & close to first-order logic; description logics & not as powerful as
  first-order logic \\
Query standards & no consolidated standard, few proposals (e.g.\ SPARQL) & consolidated
  standards (e.g.\ SQL) \\
Management systems & Protégé, Jena (not commercial) & many solid commercial tools and
  environments (Oracle, \ldots) \\
\bottomrule
\end{tabularx}
\captionof{table}{Ontologies and databases compared.}
\end{center}

\begin{nbwarning}[SPARQL today]
SPARQL has since become a W3C Recommendation (2008; SPARQL 1.1 in 2013), so it is now the
standard query language for RDF data, even if it is far less widespread than SQL. The
comparison in the slides reflects an earlier state of the field.
\end{nbwarning}

\subsubsection{Knowledge graphs}
\slidesrc{Ontologies: 19--21}

\begin{nbdef}{Knowledge graph}
``The heart of the \textbf{knowledge graph} is a knowledge model: a collection of
interlinked descriptions of concepts, entities, relationships and events. Knowledge
graphs put data in context via linking and semantic metadata and this way provide a
framework for data integration, unification, analytics and sharing'' (Ontotext).
\end{nbdef}

\begin{center}
\begin{tikzpicture}[
    font=\scriptsize,
    ent/.style={circle, draw=none, fill=nbteal!80, text=white, minimum size=1.25cm,
                align=center, font=\scriptsize\bfseries},
    co/.style={ent, fill=nbgray!80},
    typ/.style={ent, fill=nborange!90},
    rel/.style={-{Stealth}, thick, nbgray},
    inf/.style={-{Stealth}, thick, dashed, nbred!80},
    lab/.style={midway, fill=white, inner sep=1pt, font=\tiny}]
  \node[co]  (big)  at (0,  4.4) {Big\\Bucks\\Café};
  \node[co]  (glob) at (0,  2.2) {Global\\Invest.\\Inc.};
  \node[co]  (my)   at (0,  0)   {My\\Local\\Café};
  \node[ent] (sea)  at (3.4,4.4) {Seattle};
  \node[ent] (wb)   at (3.4,2.2) {West\\Bay};
  \node[ent] (nyc)  at (3.4,0)   {New York\\City};
  \node[ent] (wa)   at (6.8,4.4) {Washing-\\ton};
  \node[ent] (cay)  at (6.8,2.2) {Cayman\\Island};
  \node[ent] (ny)   at (6.8,0)   {New\\York};
  \node[ent] (usa)  at (10.2,2.2) {USA};
  \node[typ] (ctry) at (10.2,4.4) {Country};
  \node[typ] (off)  at (10.2,0)   {Offshore\\Zone};
  \draw[rel] (big)  -- node[lab]{controls}     (glob);
  \draw[rel] (glob) -- node[lab]{controls}     (my);
  \draw[rel] (big)  -- node[lab]{located in}   (sea);
  \draw[rel] (glob) -- node[lab]{located in}   (wb);
  \draw[rel] (my)   -- node[lab]{located in}   (nyc);
  \draw[rel] (sea)  -- node[lab]{subregion of} (wa);
  \draw[rel] (wb)   -- node[lab]{subregion of} (cay);
  \draw[rel] (nyc)  -- node[lab]{subregion of} (ny);
  \draw[rel] (wa)   -- node[lab, pos=0.35]{subregion of} (usa);
  \draw[rel] (ny)   -- node[lab, pos=0.22]{subregion of} (usa);
  \draw[rel] (usa)  -- node[lab]{type}         (ctry);
  \draw[rel] (cay)  -- node[lab, pos=0.8]{type} (off);
  \draw[inf] (big.north) to[bend left=25] node[lab]{located in} (wa.north);
  \draw[inf] (my.south)  to[bend right=25] node[lab]{located in} (ny.south);
  \draw[inf] (glob) to[bend left=22] node[lab]{located in} (cay);
\end{tikzpicture}
\captionof{figure}{A knowledge graph of companies and places, after the slides. Solid
arrows are stated facts, dashed arrows are \emph{inferred} facts.}
\label{fig:knowledge-graph}
\end{center}

\begin{nbexample}{what the graph lets us infer}
The stated facts say that \emph{My Local Café} is located in \emph{New York City}, which
is a subregion of \emph{New York}: since \emph{located in} propagates along
\emph{subregion of}, the café is also located in New York, and in the USA. Following the
\emph{controls} links upwards, \emph{My Local Café} is controlled by \emph{Global
Investment Inc.}, located in West Bay, a subregion of the Cayman Islands, which are an
\emph{Offshore Zone}. The slide's conclusion, ``discover suspicious relationships'', is
a query over the graph: \emph{a company controlled by another company from the same
country through an offshore zone}. A news item saying that ``My Local Café is a network
of independent coffee shops'' can be checked against it.
\end{nbexample}

\begin{nbintuition}[ontology vs.\ knowledge graph]
The two terms overlap heavily. A useful way to keep them apart: the \emph{ontology} is
the schema (classes such as \emph{Company}, \emph{Region}; properties such as
\emph{controls}, \emph{located in}; rules such as ``\emph{located in} is transitive
through \emph{subregion of}''), while the \emph{knowledge graph} is the large network of
concrete entities and facts described with that schema.
\end{nbintuition}
